%% Cover Letter - Laith Assaf
%% Built on moderncv template

\documentclass[11pt,letterpaper,sans]{moderncv}

\moderncvcolor{blue}
\moderncvstyle{classic}

\usepackage[scale=0.75]{geometry}
\setlength{\footskip}{136.00005pt}

% font loading
\ifxetexorluatex
  \usepackage{fontspec}
  \usepackage{unicode-math}
  \defaultfontfeatures{Ligatures=TeX}
  \setmainfont{Latin Modern Roman}
  \setsansfont{Latin Modern Sans}
  \setmonofont{Latin Modern Mono}
  \setmathfont{Latin Modern Math}
\else
  \usepackage[T1]{fontenc}
  \usepackage{lmodern}
\fi

\usepackage[english]{babel}

% personal data
\name{Laith}{Assaf}
\address{East Lansing, MI}{}{}
\phone[mobile]{(248) 480-5783}
\email{assaflai@msu.edu}
\social[linkedin]{laith-assaf-}
\social[github]{Novapool}

%----------------------------------------------------------------------------------
%            content
%----------------------------------------------------------------------------------
\begin{document}

%-----       letter       ---------------------------------------------------------
\recipient{ICER Recruiting Team}{Institute for Cyber-Enabled Research\\Michigan State University}
\date{\today}
\opening{Dear ICER Team,}
\closing{Best regards,}
\makelettertitle

Helping researchers spend less time wrestling with job schedulers and more time doing science is a problem I find genuinely interesting. I am a junior Computer Science student at Michigan State University applying for the \textbf{LLM Research and Development Technician} position at ICER. Building an \textbf{AI agent} that understands a user's code and translates it into a correct HPCC submission is exactly the kind of practical, measurable AI work I want to be doing.

My background maps directly to what this role requires. I am fluent in \textbf{Python} and comfortable on the command line, having administered Linux servers and Docker infrastructure for years on my personal home lab. Through MSU's AI Club, I lead workshops for 500+ students and built a production application integrating a \textbf{Python} backend with the OpenAI API, giving me hands-on experience with \textbf{agentic AI} pipelines. At MHacks 2025, I built SoundSense, a real-time audio classification system with a \textbf{machine learning} backend, winning 2nd place, which required designing validation tests to confirm model accuracy under production conditions. I also built Nexus, a server management platform in Python, where I developed testing and monitoring pipelines to evaluate system reliability across concurrent operations. Designing evaluation schemas to measure how well an agent interprets and transforms user code is a natural extension of that work.

ICER's position is unique because it combines agent development with real research outcomes: the results I gather will be presented to MSU leadership and at research forums. That accountability to a concrete research question, not just shipping a feature, is the kind of environment where I do my best work. I would welcome the chance to contribute to this project.

\makeletterclosing

\end{document}
